\documentclass[../../main.tex]{subfiles}
\begin{document}
\begin{oframed}
  \subfile{problem}
\end{oframed}

{\bf [解答]}

\paragraph{(1)}

$A=
\begin{pmatrix}a&b\\c&d
\end{pmatrix}$，$B=
\begin{pmatrix}e&f\\g&h
\end{pmatrix}$とすると
\begin{align}
  AB=
  \begin{pmatrix}ae+bg&af+bh\\ce+dg&cf+dh
  \end{pmatrix}
\end{align}
より
\begin{align}
  \Delta(AB)&=(ae+bg)(cf+dh)-(af+bh)(ce+dg)\\
  &=acef+adeh+bcfg+bdgh-acef-adfg-bceh-bdgh\\
  &=ad(eh-fg)-bc(eh-fg) \\
  &=(ad-bc)(eh-fg)=\Delta(A)\Delta(B) \label{eq:1}
\end{align}
より題意は示された．

\paragraph{(2)}

$x=\Delta(A)$，$y=t(A)$とする．題意の条件
\begin{align}
  A^5 = E \label{eq:2}
\end{align}
について考える．両辺に$\Delta$を作用させると
\begin{align}
  \Delta(A^5) = \Delta(E) = 1 \label{eq:3}
\end{align}
である．

一方\cref{eq:1}に代入して$B=A$をくりかえし用いると
\begin{align}
  \Delta(A^5)=\Delta(A)^5=x^5 \label{eq:4}
\end{align}
である．\cref{eq:3,eq:4}より
\begin{align}
  x^5 = 1
\end{align}
である．$A$の成分が全て実数だから$x$も実数なので，
\begin{align}
  x = 1 \label{eq:5}
\end{align}
が従う．

そこで以下$y$を求める．
ケイリー・ハミルトンの定理より
\begin{align}
  A^2=yA-xE = yA-E \label{eq:6}
\end{align}
だから，$A^3$に繰り返し\cref{eq:6}を用いると
\begin{align}
  A^3 = A(yA-E) = yA^2-A = y(yA-E)-A = (y^2-1)A-yE
\end{align}
より$A^5$は
\begin{align}
  A^5
  & = A^2\cdot A^3 \\
  & = (yA-E)\left\{(y^2-1)A-yE\right\} \\
  & = y(y^2-1)A^2 - (2y^2-1)A + yE \\
  & = y(y^2-1)(yA-E) - (2y^2-1)A + yE \\
  & = (y^4-3y^2+1)A + y(-y^2+2)E
\end{align}
となる．これが$E$に等しいから
\begin{align}
  (y^4-3y^2+1)A = (y^3-2y+1)E \\
  (y^2+y-1)(y^2-y-1)A = (y-1)(y^2+y-1)E \label{eq:7}
\end{align}
である．

まず$y^2+y-1=0$の時は\cref{eq:7}は成立するから，
\begin{align}
  y = \dfrac{-1\pm \sqrt{5}}{2} \label{eq:8}
\end{align}
は解である．この時\cref{eq:5}より$0\le\theta<2\pi$に対して
\begin{align}
  A =
  \begin{pmatrix}\cos\theta&\sin\theta\\ -\sin\theta&\cos\theta
  \end{pmatrix}
\end{align}
と書けるから，
\begin{align}
  y = 2\cos\theta = \dfrac{-1\pm \sqrt{5}}{2}
  \quad\therefore
  \cos\theta = \dfrac{-1\pm \sqrt{5}}{4} = \cos\dfrac{2\pi}{5}, \cos\dfrac{4\pi}{5}
\end{align}
となってこの場合の$A$は存在する．

それ以外の時，\cref{eq:7}の両辺を$y^2+y-1 \neq 0$で割って
\begin{align}
  (y^2-y-1)A = (y-1)E
\end{align}
を得る．$y^2-y-1=0$つまり$y=\dfrac{1\pm\sqrt{5}}{2}$の時は左辺は$0$だが右辺が$0$とならずに不適だから，$y^2-y-1\neq 0$で両辺を割って，
\begin{align}
  A = \dfrac{y-1}{y^2-y-1}E \label{eq:9}
\end{align}
となり$A$はスカラー行列となるから
\begin{align}
  A = \dfrac{t(A)}{2}E = \dfrac{y}{2}E
\end{align}
と書ける．\cref{eq:2}に代入して$y$が実数であることに注意して
\begin{align}
  \left(\dfrac{y}{2}\right)^5 = 1 \quad \therefore y = 2 \label{eq:10}
\end{align}
となる．この時\cref{eq:9}も満たされて十分である．
以上\cref{eq:5,eq:8,eq:10}から
\begin{align}
  x = 1, \quad
  y = 2, \dfrac{-1\pm\sqrt{5}}{2}
\end{align}
である．

\newpage

\end{document}
